\chapter{Conclusions}
\label{cha:conclusions}

This thesis set out to measure how much two grounding and review mechanisms -- a deterministic nutrient totaller and a critique-and-refine loop -- contribute to the reliability of an LLM-based one-day meal planner, on top of a food knowledge base the planner is already required to query rather than guess from. The resulting system enforces that requirement structurally: its language model never has a field in which to write a nutrient value, so every number in a delivered plan is traceable to a real food database entry and an exact arithmetic computation (\Cref{sec:impl-portion-ref,sec:design-no-invented-foods}).

Against that grounded baseline, H1 is supported through the cheaper of the two mechanisms: the totaller cuts energy and macronutrient error sharply, does so in every scenario, and makes plans more predictable. H2 is unsupported: the critique-and-refine loop adds nothing measurable on top of the totaller on either evaluation dimension, and it costs roughly three times the wall-clock time. H3 is supported in all three of its parts -- plan quality varies with the model and with reasoning effort, and raising a weaker model's effort recovers a substantial part of the gap -- with the practical caveat that spending the same latency on a stronger model at low effort still dominates. The numbers behind these verdicts are in \Cref{sec:results-ablation,sec:results-model-comparison,sec:results-overhead}.

Two secondary findings are worth carrying forward. The two evaluation dimensions do not order the configurations the same way, so neither is a proxy for the other: the two configurations tied for lowest energy and macronutrient error differ sharply in how well the judges think their plans answer the request, and reasoning effort improves accuracy monotonically while leaving preference satisfaction flat past its middle setting. And the two judges, while agreeing on almost every ranking, disagree systematically on absolute scores and, less excusably for a question with an objectively checkable answer, on where allergen and diet-pattern compliance actually slipped, which suggests hard exclusions deserve a deterministic checker rather than a judged score.

Those verdicts should be read inside the bounds of \Cref{sec:results-threats}: a five-scenario grid, LLM judges without a human expert, vendor-defined effort labels, a single non-interactive session, and a snapshot of hosted models that will not stay still.

\section{Further work}
\label{sec:conclusions-future-work}

Several directions follow directly from the limitations of \Cref{sec:results-threats}.

\textbf{Retrieval-augmented clinical guidance.} The dietary guardrails in every agent's prompt are stated inline rather than retrieved, because they apply near-universally to a healthy adult (\Cref{sec:design-agents}). Guidance that varies substantially with a diagnosed condition does not fit that reasoning, and the chronic-condition scenarios are exactly where it would matter most. Retrieving condition-specific guidance from a curated clinical corpus, in the manner \textsc{DietQA} applies to recipe data (\Cref{sec:background-related-systems}), would add a third grounding mechanism to the two measured here -- and, being independently switchable like the others, one whose contribution the same ablation design could quantify.

\textbf{Probability-weighted judge scores.} A judge that exposed token probabilities would yield G-Eval's weighted score directly (\Cref{sec:methodology-qualitative}), removing both the sampling cost of the repetition axis and the residual judge noise that \Cref{sec:results-stability} has to report around it.

\textbf{Pairwise comparison between configurations.} Every qualitative score here is absolute: a plan is read against a rubric without reference to any other plan. The judge literature's stronger signal is a forced choice between two plans for the same profile, which agrees with human preference on \qty{85}{\percent} of non-tied comparisons (\Cref{sec:background-evaluation})~\cite{zheng2023judging}. Adding it would sharpen precisely the comparisons this thesis found hardest to call -- the configurations that tied on accuracy but not on preference -- at the cost of a comparison count quadratic in the number of variants and of randomising presentation order against the documented position bias.

\textbf{Asymmetric critic model.} Whether a critic benefits from being a stronger model than the generator it reviews was left open in \Cref{sec:impl-critic-design}. Since the critique-and-refine loop's null result (\Cref{sec:results-ablation}) was obtained under an equal-capability critic, it bounds only what a self-review by an equally capable model achieves.

\textbf{Multi-turn interactive refinement.} Extending the system to accept a follow-up request against an already-generated plan, rather than only a fresh session each time, was excluded from the single-session scope set in \Cref{sec:vision-usage-scenario} as orthogonal to this thesis's research question; it remains a natural and comparatively well-scoped extension, since the critique-and-refine loop's revise-and-recheck pattern (\Cref{sec:design-reflection}) already provides most of the machinery a human-in-the-loop version of the same pattern would need.

\textbf{Model-generated food descriptions for retrieval.} The text embedded for dense retrieval is assembled by a fixed field-concatenation rule (\Cref{sec:impl-hybrid-search}). Having a language model instead author a richer description of each food catalogue entry could normalise inconsistent crowd-sourced source data more gracefully, at the cost of a second, offline LLM-dependent build step; measuring whether such descriptions actually improve retrieval quality, and by how much relative to that added build cost, is left as an open question.

\textbf{Elicitation of missing profile data.} The system currently assumes a complete, already-known user profile; a natural extension is a short elicitation dialogue that asks a small number of targeted follow-up questions when a profile is incomplete in a way that would materially affect the computed nutrient targets, rather than requiring every field to be supplied up front.
